% =============================================================================
% 8. CONCLUSION
% =============================================================================
\section{Conclusion}
\vspace{-0.5ex}
\label{sec:conclusion}

We started from the observation that a single step of a coding agent and
a single function call site share the same four-part structure---\textit{context,
action, externally produced return, continuation}---making source code an
internet-scale supply of agent-relevant signal. We turn this into
function-aware FIM mid-training: a self-supervised stage that masks
targets selected via program dependency graph analysis and a
complexity--inferability double criterion, with chain-of-thought
rationales embedded inside the FIM middle. Across three axes of
robustness---model size (7B/14B/32B), post-training pipeline (r2e-gym,
SWE-Smith, swe-lego), and base-model family (Qwen2.5-Coder, Qwen3)---the
recipe delivers consistent gains on coding-agent benchmarks, and the
same Python-only corpus also transfers to non-coding tool-use benchmarks
($\tau$-bench, BFCL) where it has no a priori reason to help. Future
work includes extending function-aware selection to non-Python languages
and composing mid-training with reinforcement-learning post-training.
